\subjecttitle{WIRELESS COMMUNICATION}{EX 715}{3}{0}{0}{IV}{I}

\cobjectives
  To introduce the student to the principles and building blocks of wireless communications.
  \VS

\begin{enumerate}
  \item \textbf{Introduction \hfill (2 hours)}
    \begin{enumerate}
      \item Evolution of wireless (mobile) communications, worldwide market, examples
      \item Comparison of available wireless systems, trends
      \item Trends in cellular radio (2G, 2.5G, 3G, beyond 3G) and personal wireless communication systems
    \end{enumerate}
    \vs

  \item \textbf{Cellular mobile communication concept \hfill (4 hours)}
    \begin{enumerate}
      \item Frequency re-use and channel assignment strategies
      \item Handoff strategies, types, priorities, practical considerations
      \item Interference and system capacity, co-channel and adjacent channel interference, power control measures
      \item Grade of service, definition, standards
      \item Coverage and capacity enhancement in cellular network, cell splitting, sectoring, repeaters, microcells
    \end{enumerate}
    \vs

  \item \textbf{Radio wave propagation in mobile network environment \hfill (12 hours)}
    \begin{enumerate}
      \item Review Free space propagation model, radiated power and electric field
      \item Review Propagation mechanisms (large-scale path loss) - Reflection, ground reflection, 
      \item Practical link budget design using path loss models.
      \item Outdoor propagation models (Longley-Rice, Okumura, Hata, Walfisch and Bertoni, microcell)
      \item Indoor  propagation models (partition losses, long-distance path loss, multiple breakpoint, attenuation factor)
      \item Small scale fading and multipath (factors, Doppler shift), Impulse response model of multipath channel, multipath measurements, parameters of mobile multipath channel (time dispersion, coherence bandwidth, Doppler spread and coherence time)
      \item Types of small-scale fading (flat, frequency selective, fast, slow), Rayleigh and Ricean fading distribution
    \end{enumerate}
    \vs

  \item \textbf{Modulation-Demodulation methods in mobile communications \hfill (4 hours)}
    \begin{enumerate}
      \item Review of amplitude (DSB, SSB, VSB) and angle (frequency, phase) modulations and demodulation techniques
      \item Review of line coding,  digital linear   (BPSK, DPSK, QPSKs) and constant envelop (BFSK, MSK, GMSK) modulation and demodulation techniques
      \item M-ary (MPSK, MFSK, QAM and OFDM) modulation and demodulation techniques
      \item Spread spectrum modulation techniques, PN sequences, direct sequence and frequency hopped spread spectrums
      \item Performance comparison of modulations techniques in various fading channels
    \end{enumerate}
    \vs

  \item \textbf{Equalization and diversity techniques \hfill (4 hours)}
    \begin{enumerate}
      \item Basics of equalization. Equalization in communications receivers, linear equalizers
      \item Non-linear equalization, decision feedback and maximum likelihood sequence estimation equalizations
      \item Adaptive equalization algorithms, zero forcing, least mean square, recursive least squares algorithms, fractionally spaced equalizers
      \item Diversity methods, advantages of diversity, basic definitions
      \item Space diversity, reception methods (selection, feedback, maximum ratio and equal gain diversity)
      \item Polarization, frequency and time diversity
      \item RAKE receivers and interleaving
    \end{enumerate}
    \vs
    
  \item \textbf{Speech and channel coding fundamentals \hfill (4 hours)}
    \begin{enumerate}
      \item Characteristics of speech signals, frequency domain coding of speech (sub-band and adaptive transform coding)
      \item Vocoders (channel, formant, cepstrum and voice-excited), Linear predictive coders (multipulse, code and residual excited LPCs), Codec for GSM mobile standard
      \item Review of block codes, Hamming, Hadamard, Golay, Cyclic, Bosh-Chaudhary- Hocquenghgem (BCH), Reed-Solomon (RS) codes
      \item Convolutional codes, encoders, coding gain, decoding algorithms (Viterbi and others)
      \item Trellis Code Modulation (TCM), Turbo codes
    \end{enumerate}
    \vs
    
  \item \textbf{Multiple Access in Wireless communications \hfill (9 hours)}
    \begin{enumerate}
      \item Frequency Division Multiple Access (FDMA), principles and  applications
      \item Time Division Multiple Access (TDMA), principles and applications
      \item Spread Spectrum Multiple Access, Frequency Hopped Multiple Access, Code Division
      \item Space Division Multiple Access
      \item Standards for Wireless Local Area Networks
    \end{enumerate}
    \vs

  \item \textbf{Wireless systems and standards \hfill (6 hours)}
    \begin{enumerate}
      \item Evolution of wireless telephone systems: AMPS, PHS, DECT, CT2, IS-94, PACS, IS-95, IS-136, IS-54 etc.
      \item Global system for Mobile (GSM): Services and features, system architecture, radio sub-system, channel types (traffic and control), frame structure, signal processing, example of a GSM call
      \item CDMA standards: Frequency and channel specifications, Forward and Reverse CDMA channels
      \item WiFi, WiMAX, UMB, UMTS, CDMA-EVDO, LTE, and recent trends
      \item Regulatory issues (spectrum allocation, spectrum pricing, licensing, tariff  regulation and interconnection issues)
    \end{enumerate}
    \vs
\end{enumerate}
\Vs

\noindent \textbf{Practical:}
  \begin{enumerate}[label=\arabic*.]
    \item Case Study and Field Visit
    \item Visits to mobile service operators, network service providers, internet service providers
  \end{enumerate}
\Vs

\noindent \textbf{References:}
  \begin{enumerate}[label=\arabic*.]
    \item K. Feher,  Wireless Digital Communications, latest editions
    \item T. Rappaport, Wireless Communications, Latest editions
    \item J. Schiller, Mobile Communications
    \item Leon Couch, Digital and analog communication systems, latest edition
    \item B.P.Lathi, Analog  and Digital communication systems, latest edition
    \item J. Proakis, Digital communication systems, latest edition
    \item D. Sharma, Course manual “Communication Systems II”.  
  \end{enumerate}

\newpage

